\subsection{Extensions of ordered fields}

DEFINITION 3. --- Let K be an ordered field. An ordered extension of K is a pair (E, u), where E is an ordered field and u is an increasing homomorphism from K to E.

Let K be a field, let E be an ordered field and let \(u: K \rightarrow E\) be a homomorphism. The relation

\[
x \leqslant y \quad \text { if } \quad u (x) \leqslant u (y)
\]

is a total order relation on K which gives it an ordered field structure, said to be induced by that of E. If K and E are ordered fields, then a homomorphism \(u: K \rightarrow E\) is increasing if and only if the ordered field structure of K is induced by that of E. We will usually identify K with its image in E under u.

Examples. --- 1) Every ordered field K is an ordered extension of Q. Indeed K is an extension of Q, since it is of characteristic zero, and on the other hand Q can only be ordered in one way, as we have just seen.

\begin{enumerate}
\def\labelenumi{\arabic{enumi})}
\setcounter{enumi}{1}
\tightlist
\item
  Let K be an ordered field, and let K(X) be the field of rational functions in one indeterminate over K. Let us define an ordering on the polynomial ring K{[}X{]} by taking the positive elements to be 0 and those polynomials whose leading coefficient is positive. In this way we obtain a totally ordered ring whose ordering extends that of K. By applying Prop. 2 we give K(X) the structure of an ordered extension of K. * For K = R it can be shown that the order relation defined on K(X) in this way is that of growth near +∞ (cf.~VI, p.~24, Prop. 4). *
\end{enumerate}

THEOREM 1. --- For an extension E of K to admit the structure of an ordered extension of K, the following condition is necessary and sufficient :

(OE) The relation \(p_{1}x_{1}^{2} + \cdots + p_{n}x_{n}^{2} = 0\) implies

\[
p _ {1} x _ {1} = \dots = p _ {n} x _ {n} = 0
\]

for any finite sequence \((x_{i},p_{i})\) of pairs of elements \(x_{i}\) of E and positive elements \(p_i\) of K.

Condition (OE) is clearly equivalent to :

(OE') The element \(-1\) is not a sum of elements of the form \(px^2\) ( \(x \in \mathbf{E}\) , \(p \in \mathbf{K}\) , \(p \geqslant 0\) ).

Condition (OE) is necessary: if E is an ordered extension of K then the elements \(p_{i}x_{i}^{2}\) are positive in E, so zero if their sum is zero. On the other hand \(p_{i}x_{i}^{2}=0\) is equivalent to \(p_{i}x_{i}=0\) .

Conversely, suppose condition (OE) is satisfied, then we will define an ordering on E by constructing a subset P of E which satisfies conditions \((\mathrm{AP}_{\mathrm{I}})\) , \((\mathrm{AP}_{\mathrm{II}})\) , \((\mathrm{AP}_{\mathrm{III}})\) and \((\mathrm{AP}_{\mathrm{IV}})\) , and which contains the set \(K_{+}\) of positive elements of K. Such a subset P will certainly make E an ordered extension of K, for we will have \(K \cap P = K_{+}\) ; indeed, if P were to contain an element -a \textless{} 0 of K, then a would belong to \(P \cap (-P)\) , contradicting \((\mathrm{AP}_{\mathrm{III}})\) .

To define P, let us consider the set M of subsets of E which satisfy \((\mathrm{AP}_{\mathrm{I}})\) , \((\mathrm{AP}_{\mathrm{II}})\) and \((\mathrm{AP}_{\mathrm{III}})\) , and which contain the union of \(K_{+}\) and the set C of squares of elements of E. This set M is nonempty, for it contains the set \(P_{0}\) of elements of the form \(\sum_{i} p_{i} x_{i}^{2}\) (that \(P_{0}\) satisfies \((\mathrm{AP}_{\mathrm{III}})\) follows immediately

from (OE)).

Moreover M is inductive (Set Theory, III, p.~154, Def. 3). Thus there exists, by Th. 2 of Set Theory, III, p.~154, a maximal element in M, which it remains for us to prove satisfies \((\mathrm{AP}_{\mathrm{IV}})\) ; now this follows from the following lemma:

Lemma. --- Let \(P \in M\) and \(x \notin P\) ; then there exists \(P' \in M\) such that \(P \subset P'\) and \(-x \in P'\) .

Take \(P' = P - xP\) , and check that \(P'\) has the required properties. Since \(0 \in C \subset P\) , we have \(P \subset P'\) . Whence \(C \subset P'\) and \(K_{+} \subset P'\) . Since \(1 \in C \subset P\) we have \(-x \in P'\) . We have

\[
\mathrm{P} ^ {\prime} + \mathrm{P} ^ {\prime} = \mathrm{P} - x \mathrm{P} + \mathrm{P} - x \mathrm{P} = \mathrm{P} + \mathrm{P} - x (\mathrm{P} + \mathrm{P}) \subset \mathrm{P} - x \mathrm{P} = \mathrm{P} ^ {\prime},
\]

whence \((\mathbf{AP}_1)\) . We have

\[
\begin{array}{r l} \mathrm {P^ {\prime} P^ {\prime}} = & (\mathrm{P} - x \mathrm{P}) (\mathrm{P} - x \mathrm{P}) \subset \\ & \subset \mathrm{PP} + x ^ {2} \mathrm{PP} - x (\mathrm{PP} + \mathrm{PP}) \subset \mathrm{P} + \mathrm{CP} - x \mathrm{P} \subset \mathrm{P} - x \mathrm{P} = \mathrm {P^ {\prime}}, \end{array}
\]

whence \((\mathrm{AP}_{\mathrm{II}})\) . Finally, let us check \((\mathrm{AP}_{\mathrm{III}})\) : suppose given an identity of the form \(p - xq = -(r - xs)\) where p, q, r, s belong to P; we deduce from this the relation \(x(s + q) = p + r\) ; if \(s + q \neq 0\) we have

\[
x = (s + q) ^ {- 2} (s + q) (p + r) \in \mathrm{CPP} \subset \mathrm{P},
\]

contrary to the hypothesis; hence \(s + q = 0\) , whence \(p + r = 0\) ; since \(\mathbf{P}\) satisfies \((\mathbf{AP}_{\mathrm{III}})\) we deduce that \(s = q = r = p = 0\) , which completes the proof.

COROLLARY 1 (« Artin-Schreier Theorem »). --- A necessary and sufficient condition for there to exist an ordering on a commutative field E which makes it an ordered field, is that the relation \(x_{1}^{2} + \cdots + x_{n}^{2} = 0\) imply \(x_{1} = \cdots = x_{n} = 0\) .

The necessity is obvious. Conversely, the stated condition implies that E is of characteristic zero, hence an extension of Q; then condition (OE) is satisfied, and Th. 1 shows that there exists on E the structure of an ordered extension of Q, that is an ordered field structure.

There does not exist any ordered field structure on a field E in which -1 is a square, in particular on an algebraically closed field.

COROLLARY 2. --- Let E be an extension of K admitting the structure of an ordered extension of K. For an element \(x \in E\) to be positive under every such structure on E, it is necessary and sufficient that x be of the form \(\sum_{i} p_{i} x_{i}^{2}\) , where \(x_{i} \in E\) and the

\(p_i\) are positive elements of \(\mathbf{K}\) .

The condition is obviously sufficient; it is also necessary, for (in the notation of the proof of Th. 1), if \(x \notin \mathbf{P}_0\) there exists a maximal element \(\mathbf{P}\) of \(\mathfrak{M}\) such that \(x \notin \mathbf{P}\) ; then \(-x \in \mathbf{P}\) by the Lemma, and \(x\) is not positive under the ordering defined by \(\mathbf{P}\) , since \(x \neq 0\) .

